\subsection{Typed transitions and policy-relative satisfaction}
\label{app:typed-contract}
The paper uses P2 for context/evidence integrity and P3 for principal authority
and the authorized value throughout. Cross-candidate context agreement is a
separate binding condition; it is not the internal integrity condition P2.

Let $\Sigma$ contain the current version, complete field-value and field-source
maps, and the candidate, authorization and transition registries. A submitted
item is $u=(c_s,a,x)$, where $a=(c_r,x_a,q)$. For the singleton presentation, a
persisted transition has type
\[
\begin{aligned}
 \tau&=\langle j,r,f,e,v^-,v^+,x,p,\operatorname{id}(c_s),\operatorname{id}(a)\rangle,\\
 \tau&\in \mathsf{Tid}\times\mathsf{Record}\times\mathsf{Field}\times\mathsf{Evidence}
 \times\mathsf{Version}^{2}\\
 &\qquad\times\mathsf{Value}\times\mathsf{Producer}
 \times\mathsf{Cid}\times\mathsf{Aid}.
\end{aligned}
\]
The transition identifies the executed candidate; the authorization separately
identifies the reviewed candidate; $\mathsf{Cid}$ and $\mathsf{Aid}$ are registry
keys. Certificate rows in the realization resolve
these references to immutable candidate records. The fields of this tuple are
observable relations, not an assumption that an arbitrary stored tuple is valid.

For $b\in\{\mathrm{context},\mathrm{bound}\}$, define the binding condition
\[
\begin{aligned}
 \beta_{\mathrm{context}}(u)&\iff K(c_s)=K(c_r),\\
 \beta_{\mathrm{bound}}(u)&\iff\beta_{\mathrm{context}}(u)
       \land\operatorname{id}(c_s)=\operatorname{id}(c_r).
\end{aligned}
\]
Thus P1 denotes $\beta_b$ under the selected policy. Under bound admission it
includes instance equality; under context admission it requires only the
declared context equality. In particular, context admission does not borrow
instance equality through another shared check.

Write $z=(\Sigma^-,u,\tau,\Sigma^+)$ for a committed outcome. The satisfaction
relation is defined by the following checks on the outcome and the validated
admission inputs:
\begin{description}
\item[P1$_b$:] $\beta_b(u)$ holds for the submitted and reviewed records.
\item[P2:] The declared record, field, evidence identity/locator and expected
predecessor agree with the submitted candidate. The evidence belongs to the
target record; the transition retains the candidate's record, field, evidence
and producer. This concerns internal consistency, independently of P1$_b$.
\item[P3:] The authorization comes from a permitted principal at commit entry,
is usable for this admission, and explicitly records $x_a$. The submitted,
transition and newly installed field values equal $x_a$ under the declared
comparator; no equality between $x_a$ and the retained proposal $x_c$ is required.
\item[P4:] The candidate's expected predecessor equals the authoritative
version observed in $\Sigma^-$ within the admission transaction.
\item[P5:] Exactly one successor of that predecessor is installed. Its changed
fields and new transitions correspond exactly to the admitted items, each
transition records the actual predecessor and successor, and all other values
are retained. The successor, new transitions, source updates and grant
consumption belong to the same commit unit.
\item[P6$_{\mathrm{local}}$:] Each changed field points to its unique new
transition, with the same field and installed value. Each unchanged field
retains its exact prior source. Source references resolve uniquely in the
retained transition registry.
\end{description}
For a batch, $u$ and $\tau$ are replaced by the admitted item set and its
transition set; P1--P4 apply itemwise and P5 requires one shared record/version
envelope and an item--transition bijection. After initial coverage, the concrete
admitted set contains exactly the submitted fields whose canonical values
change and must be nonempty, as specified in Online Resource 1, S2.

Formally, with $L(z)$ denoting P2--P5 and P6$_{\mathrm{local}}$ together,
\[
 z\models\mathcal R_b^{\mathrm{local}}
 \quad\Longleftrightarrow\quad \beta_b(u)\land L(z).
\]
P0 restricts the execution model: candidate-producing operations cannot write
authoritative values, versions, sources, transitions or authorizations. It is
not inferred from $z\models\mathcal R_b^{\mathrm{local}}$. The pre/post relation
specifies one complete step; crash atomicity cannot be inferred from endpoint
snapshots. Trusted immutable lookup, canonicalization and atomic installation
are explicit execution premises. The shared gate
$G$ implements the common checks and lifecycle constraints; its definition
must not hide the bound-policy conjunct in $\beta_b$.

Full P6 also checks the source--transition--authorization--candidate--evidence
chain, including retained references, field/value/producer/version consistency
and source-version uniqueness. Its binding link is interpreted under the
declared policy: a context-authorized transition retains both candidate roles
and requires $\beta_{\mathrm{context}}$, whereas an exact-instance trace also
requires $\beta_{\mathrm{bound}}$. A context-legal substitution can therefore
satisfy its local contract while failing the separate exact-instance trace
criterion. The positive formal preservation result below is for the bound
model only; this policy-relative definition adds no new model-checking result.

\subsection{Paper obligations and archived Alloy assertions}
\label{app:formal-mapping}
Table~\ref{tab:formal-contract-map} is a semantic crosswalk, not a renaming of
the frozen source or a claim that every paper-level premise was model checked.
\textsf{H} denotes the historical singleton model
\nolinkurl{formal/alloy/auto_decte.als}; \textsf{B} denotes the batch model
\nolinkurl{latest/formal/alloy/batch/auto_decte_batch.als} in the fixed reference
artifact cited in Online Resource 1, S1. The \textsf{B} command manifest uses
freeze \nolinkurl{2026-09-11-trace-witness}. Assertion names retain their
historical spelling and numbering: neither source has an assertion named P2.

\begin{table}[htbp]
\centering\footnotesize
\caption{Paper obligations mapped to exact archived assertion names. \textsf{H}
and \textsf{B} are source namespaces, not additional evidence counts.}
\label{tab:formal-contract-map}
\begin{tabularx}{\textwidth}{@{}>{\raggedright\arraybackslash}p{0.13\textwidth}>{\raggedright\arraybackslash}p{0.40\textwidth}>{\raggedright\arraybackslash}X@{}}
\toprule
Paper relation & Archived assertions & Encoded part and limit\\
\midrule
P0 & \textsf{H}: \nolinkurl{P0_no_machine_write}; \textsf{B}: \nolinkurl{P0NoMachineAuthoritativeWrite}
& Frame modeled machine operations. They do not prove that a deployed caller cannot bypass the gate.\\
P1$_{\mathrm{bound}}$ & \textsf{H}: \nolinkurl{P1_auth_nontransfer}; \textsf{B}: \nolinkurl{P3AuthorizationIntegrity}
& Canonical certificate equality binds authorization to the submitted certificate. The candidate abstraction must preserve this relation.\\
P1$_{\mathrm{context}}$ & No separate context-policy assertion
& Defined here through $K_{\mathrm{eq}}$. Neither Full model verifies the relaxed policy or the three-level binding comparison.\\
P2 & \textsf{H}: \nolinkurl{P1_no_cross_context_admission}; \textsf{B}: \nolinkurl{P1ExactContext}
& Item/certificate record, field, evidence and expected-version consistency. These historical ``P1'' names correspond to paper P2, not reviewed-instance equality.\\
P3 & \textsf{H}: \nolinkurl{P3_authorization_reference_integrity}; \textsf{B}: \nolinkurl{P3AuthorizationIntegrity}
& Principal membership, explicit authorized value and reference checks. Dynamic revocation, expiry and grant consumption are outside these models.\\
P4 & \textsf{H}: \nolinkurl{P4_stale_reject}; \textsf{B}: \nolinkurl{P4Freshness}
& Expected/current predecessor equality. Database concurrency semantics are an implementation obligation.\\
P5 & \textsf{H}: \nolinkurl{P5_transition_integrity}; \textsf{B}: \nolinkurl{P5TransitionBijection}, \nolinkurl{FullNoPartialItemEffects}, \nolinkurl{RejectionHasNoEffect}
& Transition relations, complete batch effects and rejection framing in atomic pre/post steps; no crash or database atomicity proof.\\
P6 & \textsf{H}: \nolinkurl{P6_trace_soundness}, \nolinkurl{P6_missing_source_incomplete}, \nolinkurl{P6_duplicate_incomplete}; \textsf{B}: \nolinkurl{P6MissingAnchorIncomplete}, \nolinkurl{P6DuplicateSourceVersionIncomplete}, and positive preservation (below)
& Missing/duplicate-source diagnostics and conditional positive trace checks; scope and pre-state restrictions remain essential.\\
\bottomrule
\end{tabularx}
\end{table}

The \textsf{B} authorization assertion calls \nolinkurl{itemAuthorizationOk},
which includes both \nolinkurl{authorizationCertificateBinding} (paper P1)
and \nolinkurl{principalBinding}/\nolinkurl{authorizationValueBinding} (paper
P3). Its context assertion calls \nolinkurl{itemContextOk} (paper P2).
This explains why the mapping is many-to-many. The Full assertions often
check consequences of explicit contract conjuncts; effective ablation
witnesses and independent persisted-state projections provide different
evidence and are not subsumed by these names.

The positive \textsf{B} P6 assertion,
\nolinkurl{LegalAdmissionTraceCompleteUnderWellFormedPre}, assumes every anchored predecessor field
has a complete trace and no prior transition targets an admitted field's next
version. It concludes trace completeness for admitted fields after one legal
Full step. The same archive includes reachable initial and carry-forward
precondition witnesses. These checks do not establish arbitrary-history
induction, every unchanged field's full chain, external-document immutability,
or the correctness of an unmodeled review channel.

For source identification, SHA-256 of \textsf{H} is
\nolinkurl{3e42226a642e458170b85671a56ff8438fdc7c8dba72b970c53822867e7f9d19}.
The \textsf{B} S1 and S2 source hashes are, respectively,
\nolinkurl{e5275f855eff19469ecb73da26d1f3dab383c2578e72f4b526a9694230f8e70b}
and \nolinkurl{3bf12ffb44259fff5022e749fe13d2573ae1539e7f8306a3402da1db0007b947}.
The matching manifest records 72 outcomes (22 SAT witnesses and 14 UNSAT
checks per profile). This crosswalk identifies the archived results; it adds
neither a new Analyzer run nor an unbounded proof.
